\begin{table}[tp]
\centering
\caption{Comparison for \href{https://github.com/k3nfalt/longtails/blob/main/LongTails/Setups/GeneralizedSmoothParameterFree.lean}{GeneralizedSmoothParameterFree}.}\label{tab:GeneralizedSmoothParameterFree}
\parbox{\linewidth}{\footnotesize Convex $f:E\to\mathbb R$ on a real Hilbert space, with differentiable gradient and $\|\nabla^2 f(x)\|\le L_0+L_1\|\nabla f(x)\|$ everywhere. A minimizer $x_\star$ satisfies $\|x_0-x_\star\|\le R$ and $f(x_0)-f(x_\star)\le\Delta$, where $L_0,R,\Delta,\varepsilon,h_0>0$ and $L_1\ge0$. Exact gradients are available. Target: $f(\hat x)-f(x_\star)\le\varepsilon$. In \textbf{(1)}, updates receive only $x_0,h_0$; the sufficient stopping budget depends on the problem parameters and $\varepsilon$.}\par\smallskip
\begin{threeparttable}
\footnotesize
\setlength{\tabcolsep}{4pt}\renewcommand{\arraystretch}{1.25}
\begin{tabularx}{\linewidth}{@{}>{\raggedright\arraybackslash}p{0.20\linewidth}>{\raggedright\arraybackslash}p{0.32\linewidth}>{\raggedright\arraybackslash}p{0.27\linewidth}>{\raggedright\arraybackslash}X@{}}
\toprule
\textbf{Result} & \textbf{Oracle complexity} & \textbf{Novelty} & \textbf{Limitations} \\
\midrule
\textbf{(1)} \href{https://github.com/k3nfalt/longtails/blob/main/LongTails/Theorems/BacktrackingGeneralizedAccelerated.lean}{Backtracking\allowbreak{}Generalized\allowbreak{}Accelerated} \textbf{(Ours)} $(\star\star\star)$ \href{https://github.com/k3nfalt/longtails/blob/main/scripts/reviews/BacktrackingGeneralizedAccelerated.md}{(review)}
& $\begin{aligned}[t]
\mathcal{O}\Bigl(&\sqrt{L_0R^2/\varepsilon}+(1+s)\log\Bigl(2+s\\
&+\frac{1}{L_0h_0}+L_0h_0\\
&+\frac{L_1^2\Delta}{L_0}+L_1^2\Delta h_0\Bigr)\Bigr)
\end{aligned}$ gradient calls, including two per trial and all rejected trials.
& & \\
\midrule[1.2pt]
\textbf{(2)} Accelerated GRAAL \citep[Cor.~3, v2]{gspfBorodichKovalev2025v2}. Restriction: $E=\mathbb R^d$, $f\in C^2$, $L_1>0$.
& $\begin{aligned}[t]
\mathcal{O}\Bigl(&1+\sqrt{L_0R^2/\varepsilon}+s^3\\
&+(1+s^2)\log\frac{1}{L_0\eta_0}\Bigr)
\end{aligned}$ gradient and value calls; requires $L_0\eta_0e^s\le1$.
& \textbf{(1)}: for $s\gg1$, $h_0=\eta_0$, $\Delta=f(x_0)-f(x_\star)$, proved start-up improves by a factor of order $s$; total gain when $\varepsilon\gtrsim L_0R^2/s^6$. Allows any $h_0>0$; removes values. Matches order for $s\lesssim1$ or accuracy-dominated costs.\par
Example: $f(x)=e^x-x$, $L_0=L_1=1$, $x_0=R\gg1$, $h_0=\eta_0=e^{-R}$: \textbf{(1)} $\mathcal{O}(R/\sqrt\varepsilon+R^2)$ versus \textbf{(2)} $\mathcal{O}(R/\sqrt\varepsilon+R^3)$.
& \textbf{(1)}: backtracking incurs constant-factor overhead. \\
\midrule
\textbf{(3)} AGD \citep[Thm.~4.3, v2]{cgsTyurin2026v2}. Restriction: $E=\mathbb R^d$; $R=\|x_0-x_\star\|$ for the displayed tuning.
& $\begin{aligned}[t]
\mathcal{O}\Bigl(&1+\sqrt{L_0R^2/\varepsilon}\\
&+s\log\Bigl(2+\min\Bigl\{\frac{L_1^2\Delta}{L_0},\frac{\Delta}{\varepsilon}\Bigr\}\Bigr)\\
&+\log\Bigl(2+\frac{\Delta}{L_0R^2}\Bigr)\Bigr)
\end{aligned}$ gradient calls; supplied $L_0,L_1$, $\bar R=R$, $\Gamma_0=2\Delta/R^2$.
& \textbf{(1)}: removes $L_0,L_1,R,\Delta$ from updates; matches the leading $\varepsilon^{-1/2}$ term. No established rate improvement.
& \textbf{(1)}: possible $(1+s)\log(L_0h_0+1/(L_0h_0))$ and $(1+s)\log(2+s)$ overhead; loses the shortcut when $\Delta/\varepsilon<L_1^2\Delta/L_0$. Backtracking constant factor. \\
\midrule
\textbf{(4)} AGMsDR \citep[Thm.~6.2, v3]{gspfVankovEtAl2025v3}. Restriction: $E=\mathbb R^d$, $f\in C^2$.
& $\begin{aligned}[t]
\mathcal{O}\Bigl((\nu+1)\Bigl[&1+\sqrt{L_0R^2/\varepsilon}\\
&+\lceil s^{2/3}\rceil\bigl(1+\log(\Delta/\varepsilon)\bigr)\Bigr]\Bigr)
\end{aligned}$ gradient/value calls; $0<\varepsilon\le\Delta=f(x_0)-f(x_\star)$. Requires $L_0,L_1$ and exact one-dimensional minimization.
& \textbf{(1)}: no $L_0,L_1$, values, or exact line search; start-up logarithm is independent of $\varepsilon$. Total rates are not comparable: $\nu$ has no parameter bound.
& \textbf{(1)} loses the $s^{2/3}$ start-up power, ignoring the unbounded search factor $\nu$. \\
\bottomrule
\end{tabularx}
\begin{tablenotes}[flushleft]\footnotesize
\item[] \textbf{Abbr:} $s=L_1R$; $h_0$ and $\eta_0$ are initial step sizes; $\bar R$ is a supplied distance bound; $\Gamma_0$ is AGD's initial potential parameter; $\nu$ bounds the gradient/value calls per exact one-dimensional search. $\log$ denotes the natural logarithm.
\end{tablenotes}
\end{threeparttable}
\end{table}
\clearpage
